\section*{अध्याय \ref{ch:g11:absorbance} --- रंग और अवशोषकता}

\begin{solution}{exo:g11:absorbance:1}
नारंगी अवशोषित करने पर वह नीला दिखता है, जो चक्र पर नारंगी के सामने का रंग है।
बैंगनी अवशोषित करने पर वह पीला दिखता है।
\end{solution}

\begin{solution}{exo:g11:absorbance:2}
हरा विलयन मुख्यतः लाल प्रकाश अवशोषित करता है, जो चक्र पर हरे के सामने का रंग है:
लगभग \qty{620}{nm} से \qty{750}{nm} तक।
\end{solution}

\begin{solution}{exo:g11:absorbance:3}
नीला रंजक: $\lambda_{\max} = \qty{629}{nm}$, $A = 0.82$। पीला रंजक:
$\lambda_{\max} = \qty{427}{nm}$, $A = 0.795$, लगभग 0.80।
\end{solution}

\begin{solution}{exo:g11:absorbance:4}
$A$ सांद्रता के समानुपाती है: तीन गुना अधिक सांद्रता से
$A = 0.90$; आधी सांद्रता से $A = 0.15$।
\end{solution}

\begin{solution}{exo:g11:absorbance:5}
रिक्त विलयन शुद्ध \omterm{def:g6:solutions-and-solubility:solute}{विलायक} से भरा एक क्यूवेट है। उससे $A = 0$ करने पर वह हट जाता
है जो क्यूवेट, \omterm{def:g6:solutions-and-solubility:solute}{विलायक} और उपकरण अवशोषित करते हैं, ताकि आगे के पाठ्यांक
केवल \omterm{def:g6:solutions-and-solubility:solute}{विलेय} के कारण हों।
\end{solution}

\begin{solution}{exo:g11:absorbance:6}
$C = 0.41 / 0.164 = \qty{2.5}{mg/L}$।
\end{solution}

\begin{solution}{exo:g11:absorbance:7}
\qty{629}{nm} पर \omterm{def:g11:absorbance:absorbance}{अवशोषकता} सबसे अधिक है, इसलिए मापन सबसे अधिक संवेदनशील है
(छोटी सांद्रताएँ भी पढ़ने लायक \omterm{def:g11:absorbance:absorbance}{अवशोषकता} देती हैं); और पट्ट के शिखर पर वक्र
सपाट है, इसलिए तरंगदैर्घ्य की छोटी त्रुटि पाठ्यांक को
लगभग नहीं बदलती। \qty{550}{nm} पर \omterm{def:g11:absorbance:absorbance}{अवशोषकता} कम है और तरंगदैर्घ्य के साथ
तेज़ी से बदलती है।
\end{solution}

\begin{solution}{exo:g11:absorbance:8}
तनु किया गया: $C = 0.48 / 0.164 = \qty{2.9}{mg/L}$; पेय:
$5 \times 2.9 = \qty{15}{mg/L}$ (अधिक परिशुद्ध रूप से $14.6$)।
\end{solution}

\begin{solution}{exo:g11:absorbance:9}
$a = A / (\ell\, C_m) = 0.795 / (1.00 \times 0.0150) = \qty{53.0}{L/(g.cm)}$;
$\varepsilon = a \times M = 53.0 \times 534.4 = \qty{2.83e4}{L/(mol.cm)}$।
\end{solution}

\begin{solution}{exo:g11:absorbance:10}
लगभग \qty{520}{nm} से ऊपर। \qty{629}{nm} पर पीला रंजक बिल्कुल
अवशोषित नहीं करता, इसलिए वहाँ की पूरी \omterm{def:g11:absorbance:absorbance}{अवशोषकता} नीले रंजक से आती है।
\end{solution}

\begin{solution}{exo:g11:absorbance:11}
दुगनी हो जाती है: $A = \varepsilon \ell c$, $\ell$ के समानुपाती है; प्रकाश
दुगने विलयन से गुज़रता है, इसलिए दुगने अवशोषक \omterm{def:g7:atoms-and-molecules:molecule}{अणुओं} से।
\end{solution}

\begin{solution}{exo:g11:absorbance:12}
$A/C$: $0.164$, $0.162$, $0.115$, $0.0675$। अनुपात केवल पहले दो के लिए
स्थिर है: लगभग \qty{10}{mg/L} से आगे (लगभग 1.6 से अधिक \omterm{def:g11:absorbance:absorbance}{अवशोषकताओं} पर)
नियम लागू नहीं रहता। केवल 5 और \qty{10}{mg/L} वाले मानकों (और उनसे कम
वालों) का उपयोग किया जा सकता है; 2.5 पढ़ने वाले अज्ञात को तनु करना होगा, उदाहरण
के लिए पाँच गुना, और फिर से मापना होगा।
\end{solution}

\begin{solution}{exo:g11:absorbance:13}
नीला रंजक, केवल \qty{629}{nm} से: $C = 0.574 / 0.164 = \qty{3.5}{mg/L}$।
पीला रंजक, \qty{427}{nm} से: $C = 0.53 / 0.0530 = \qty{10}{mg/L}$।
\end{solution}

\begin{solution}{exo:g11:absorbance:14}
उसे $0.368 / 0.164 = \qty{2.24}{mg/L}$ मिलता है। वास्तविक \omterm{def:g11:absorbance:absorbance}{अवशोषकता}
$0.368 - 0.040 = 0.328$ है, इसलिए वास्तविक सांद्रता
$0.328 / 0.164 = \qty{2.00}{mg/L}$ है। उसका उत्तर \qty{12}{\percent} अधिक है।
\end{solution}

\begin{solution}{exo:g11:absorbance:15}
प्रति दिन $10 \times 60 = \qty{600}{mg}$, यानी $600 / 20 = \qty{30}{L}$
पेय में। कोई भी दिन में तीस लीटर नहीं पीता: अकेला यह पेय कोई वास्तविक जोखिम नहीं
है, हालाँकि रंजक दूसरे खाद्य पदार्थों से भी आ सकता है।
\end{solution}

\begin{solution}{pb:g11:absorbance:1}
\textbf{1.} नारंगी-लाल, जो चक्र पर नीले के सामने का रंग है।

\textbf{2.} $\lambda_{\max} = \qty{629}{nm}$।

\textbf{3.} \qty{450}{nm} पर प्रकाश नीला है: रंजक उसे पार निकलने देता है,
और ठीक इसीलिए पेय नीला दिखता है।

\textbf{4.} \qty{629}{nm} पर।

\textbf{5.} $\qty{50.0}{mg} / \qty{0.5000}{L} = \qty{100}{mg/L}$।

\textbf{6.} \omterm{def:g10:concentration-and-dilution:dilution}{तनुकरण गुणांक} $100 / 3.0 = 33.3$, इसलिए
$V_0 = 100.0 / 33.3 = \qty{3.00}{mL}$: एक अंशांकित पिपेट (या एक
ब्यूरेट) और \qty{100.0}{mL} का आयतनमितीय फ़्लास्क।

\textbf{7.} $A/C$ = 0.168, 0.163, 0.165, 0.163, 0.165: सभी
0.164 के क़रीब, इसलिए $A \approx 0.164\, C$, मूलबिंदु से गुज़रती एक रेखा।

\textbf{8.} बिना रंजक का विलयन ($C = 0$) रिक्त विलयन है, जिसे
$A = 0$ पर रखा जाता है।

\textbf{9.} $a = \qty{0.164}{L/mg}$ प्रति सेंटीमीटर, यानी
\qty{164}{L/(g.cm)}; $\varepsilon = 164 \times 792.9 =
\qty{1.30e5}{L/(mol.cm)}$।

\textbf{10.} $F = 50.0 / 25.0 = 2$।

\textbf{11.} $C = 0.590 / 0.164 = \qty{3.60}{mg/L}$।

\textbf{12.} $2 \times 3.60 = \qty{7.20}{mg/L}$।

\textbf{13.} लगभग $2 \times 0.590 = 1.18$, जो सबसे ऊँचे मानक (0.823) से ऊपर है:
पाठ्यांक अंशांकित परास से बाहर होता, जहाँ नियम विफल हो
सकता है। तनु करने से वह मानकों के बीच आ जाता है।

\textbf{14.} $\qty{7.20}{mg/L} \times \qty{0.500}{L} = \qty{3.60}{mg}$।

\textbf{15.} $n = \num{3.60e-3} / 792.9 = \qty{4.54e-6}{mol}$।

\textbf{16.} प्रति दिन $6 \times 30 = \qty{180}{mg}$।

\textbf{17.} $180 / 3.60 = 50$ बोतलें।

\textbf{18.} एक दिन में $50 \times 0.500 = \qty{25}{L}$: पीना असंभव।
अकेले इस पेय का रंजक किसी बच्चे को सीमा के पास नहीं ला सकता,
हालाँकि कई खाद्य पदार्थों के रंजक जुड़ते जाते हैं।

\textbf{19.} $6 \times 70 = \qty{420}{mg}$, यानी $420 / 3.60 \approx
117$ बोतलें।

\textbf{20.} एक ही दिन में \qty{500}{mL} की लगभग 50 बोतलें, यानी \qty{25}{L} पेय।
\end{solution}
